% Minimal but structurally complete paper, checked in so the end-to-end
% conversion test runs everywhere — including CI, where the full external
% fixture is unavailable. Each construct below corresponds to a repair rule
% that has been observed to misfire at some point:
%
%   * a two-level header with \multirow / \multicolumn  -> tables rule
%   * a table caption forced inside a figure float       -> captions rule
%   * a wrapfigure caption pandoc drops                  -> captions rule
%   * a reference to each float kind                     -> refs rule
%   * an appendix with lettered sections                 -> headings rule
%   * an abstract in a custom environment                -> structure rule
%   * a display equation and an align block              -> math rule
\documentclass{article}
\usepackage{graphicx}
\usepackage{multirow}
\usepackage{wrapfig}

\newcommand{\metricpm}[2]{$#1 \pm #2$}
\newcommand{\metricbest}[1]{\textbf{#1}}

\title{A Tiny Paper for the \\ Converter}
\author[1]{Ada Lovelace\thanks{Equal contribution.}}
\author[2]{Alan Turing}

\begin{document}
\maketitle

\begin{abstract}
We test a converter on \textbf{bold} and \emph{emphasised} text with an
equation $E = mc^2$ inline.
\end{abstract}

\section{Introduction}\label{sec:intro}
See Section~\ref{sec:method}, Table~\ref{tab:main}, Figure~\ref{fig:one} and
Equation~\ref{eq:first}. Also \citet{example2026}.

\begin{equation}\label{eq:first}
E = mc^2
\end{equation}

\begin{align}\label{eq:second}
a &= b, \\
c &= d.
\end{align}

\section{Method}\label{sec:method}
The main comparison is in Table~\ref{tab:main}.

\begin{table}[t]
\caption{\textbf{Main results.} Scores are mean $\pm$ standard deviation.}
\label{tab:main}
\centering
\begin{tabular}{lcc}
\toprule
\multirow{2}{*}{\textbf{Method}} & \multicolumn{2}{c}{\textbf{Benchmark}} \\
\cmidrule(lr){2-3}
 & Success & Rounds \\
\midrule
Baseline & \metricpm{20.1}{2.1} & 17.7 \\
\metricbest{Ours} & \metricbest{\metricpm{24.2}{1.5}} & 18.0 \\
\bottomrule
\end{tabular}
\end{table}

\begin{figure}[t]
\captionsetup{type=table}
\caption{\textbf{Forced table caption.} Numbered by the table counter.}
\label{tab:forced}
\centering
\begin{tabular}{lc}
\toprule
Item & Value \\
\midrule
a & 1 \\
\bottomrule
\end{tabular}
\end{figure}

\begin{wrapfigure}{r}{0.4\linewidth}
\caption{A wrapfigure caption pandoc drops.}
\label{fig:one}
\end{wrapfigure}

Special symbols: $\alpha$, $\beta$, $\mathcal{L}$, and a percentage 42\%.

\appendix
\section{Proofs}\label{app:proofs}
The full derivation follows.

\subsection{A lemma}\label{app:lemma}
Statement and proof.

\bibliography{refs}

\end{document}
